\documentclass[12pt,twocolumn]{article}
\usepackage[margin=0.75in]{geometry}
\title{Two-Column Research Paper}
\author{Dr.\ Robert Chen \and Dr.\ Maria Lopez}
\date{2026}
\begin{document}
\maketitle

\begin{abstract}
This paper presents a two-column layout commonly used in academic publications.
The converter should correctly identify column boundaries and reconstruct the
reading order so that text flows from the left column to the right column on
each page, not interleaving lines from both columns.
\end{abstract}

\section{Introduction}
Academic papers frequently use two-column layouts to maximize information density
on each page. This presents a challenge for PDF extraction because the converter
must determine which text blocks belong to which column.

The spatial analysis must consider the horizontal gap between columns as a
separator. Text blocks on the left side of the gap should be read first, followed
by text blocks on the right side.

\section{Methodology}
Our approach uses a vertical split-point detector. We analyze the distribution
of text block x-coordinates and identify a gap region near the horizontal center
of the page. Blocks to the left of this gap are assigned to column 1, and blocks
to the right are assigned to column 2.

Within each column, blocks are ordered top-to-bottom. The final reading order
concatenates column 1 followed by column 2 for each page.

\section{Results}
Testing on 500 academic papers showed that our column detection achieves 97.3\%
accuracy in correctly ordering text blocks. The most common failure case involves
wide equations or figures that span both columns.

\section{Conclusion}
Two-column layout handling is essential for academic paper conversion. Our
spatial analysis approach provides robust column detection without requiring
machine learning or training data.
\end{document}